% GENERATED FILE. Edit canonical lessons, then run npm run generate:books.
% canonical-source-hash: fnv1a64:3e5d15c82e1601df
% canonical-lessons: ML-C488-oppu, ML-C488-pakarppu, ML-C488-tiriccariyalrekha, ML-C488-kalavadhi, ML-C488-kalaharanappetuka

\chapter{Papers and Signatures}
\label{ch:ml-papers-and-signatures}
\hlchaptermodality{488}

\begin{chapteropening}
\textbf{By the end of this chapter:} \emph{I can say a signature, a duplicate copy, an identity document, a validity period and to expire.}

Complete the last lesson of chapter 488: I can say a signature, a duplicate copy, an identity document, a validity period and to expire.
\end{chapteropening}

\section[\texorpdfstring{\ml{ഒപ്പ്} (oppŭ)}{ഒപ്പ് (oppŭ)}]{\ml{ഒപ്പ്} (oppŭ) — a signature}

\label{lesson:ML-C488-oppu}



\begin{quote}
Before the new one: say the Malayalam for a district collector, then the Malayalam for a nuisance.
\end{quote}

\subsection*{You'll want to know: \ml{ഒപ്പ്}}
\textbf{\ml{ഒപ്പ്}} — \emph{oppŭ} — \textquotedblleft{}a signature\textquotedblright{}.

\begin{quote}
\ml{ഫോമിന്റെ} \ml{താഴെ} \ml{നിങ്ങളുടെ} \ml{ഒപ്പ്} \ml{വേണം}.
\end{quote}

\begin{quote}
\emph{phōminṟe tāḻe niṅṅaḷuṭe oppŭ vēṇaṁ.}
\end{quote}

\begin{quote}
Your signature is needed at the bottom of the form.
\end{quote}

You met it in \ml{ഒപ്പിടുക}, to sign: to put your \ml{ഒപ്പ്}.

\begin{cousinweb}
\textbf{\ml{ഒപ്പ്}} (\emph{oppŭ}) is native Dravidian; the root means to agree or match (DEDR 924), and a signature is the mark of your agreement.
\end{cousinweb}

\begin{grammarlens}[title={What you've built}]
One more word for getting around.
\end{grammarlens}

\subsection*{Your turn}
\begin{itemize}
\raggedright
  \item \emph{Say it:} \emph{oppŭ}
  \item \emph{Say it:} \emph{oppŭ}, once more
  \item \emph{Say it:} \emph{phōminṟe tāḻe niṅṅaḷuṭe oppŭ vēṇaṁ.} — Your signature is needed at the bottom of the form.
  \item \emph{Recall:} say the Malayalam for a district collector, then the Malayalam for a nuisance, then say \emph{oppŭ} again
\end{itemize}

\begin{tcolorbox}[breakable,colback=teal!4,colframe=teal!35!black,title={Before you move on}]
What does \ml{ഒപ്പ്} mean? (\textquotedblleft{}A signature\textquotedblright{}.) Say it once more.
\end{tcolorbox}
\section[\texorpdfstring{\ml{പകർപ്പ്} (pakarppŭ)}{പകർപ്പ് (pakarppŭ)}]{\ml{പകർപ്പ്} (pakarppŭ) — a duplicate copy}

\label{lesson:ML-C488-pakarppu}



\begin{quote}
Before the new one: say the Malayalam for a nuisance, then the Malayalam for a signature.
\end{quote}

\subsection*{You'll want to know: \ml{പകർപ്പ്}}
\textbf{\ml{പകർപ്പ്}} — \emph{pakarppŭ} — \textquotedblleft{}a duplicate copy\textquotedblright{}.

\begin{quote}
\ml{പാസ്പോർട്ടിന്റെ} \ml{രണ്ട്} \ml{പകർപ്പുകൾ} \ml{കൊണ്ടുവരണം}.
\end{quote}

\begin{quote}
\emph{pāspōrṭṭinṟe raṇṭŭ pakarppukaḷ koṇṭuvaraṇaṁ.}
\end{quote}

\begin{quote}
You must bring two copies of the passport.
\end{quote}

Plural: \ml{പകർപ്പുകൾ}, \emph{pakarppukaḷ}.

The noun of \ml{പകർത്തുക}, to copy.

\begin{cousinweb}
\textbf{\ml{പകർപ്പ്}} (\emph{pakarppŭ}) is native Dravidian, the noun of \ml{പകർത്തുക}, to transfer or copy (DEDR 3803).
\end{cousinweb}

\begin{grammarlens}[title={What you've built}]
One more word for getting around.
\end{grammarlens}

\subsection*{Your turn}
\begin{itemize}
\raggedright
  \item \emph{Say it:} \emph{pakarppŭ}
  \item \emph{Say it:} \emph{pakarppŭ}, once more
  \item \emph{Say it:} \emph{pāspōrṭṭinṟe raṇṭŭ pakarppukaḷ koṇṭuvaraṇaṁ.} — You must bring two copies of the passport.
  \item \emph{Recall:} say the Malayalam for a nuisance, then the Malayalam for a signature, then say \emph{pakarppŭ} again
\end{itemize}

\begin{tcolorbox}[breakable,colback=teal!4,colframe=teal!35!black,title={Before you move on}]
What does \ml{പകർപ്പ്} mean? (\textquotedblleft{}A duplicate copy\textquotedblright{}.) Say it once more.
\end{tcolorbox}
\section[\texorpdfstring{\ml{തിരിച്ചറിയൽരേഖ} (tiriccaṟiyalrēkha)}{തിരിച്ചറിയൽരേഖ (tiriccaṟiyalrēkha)}]{\ml{തിരിച്ചറിയൽരേഖ} (tiriccaṟiyalrēkha) — an identity document}

\label{lesson:ML-C488-tiriccariyalrekha}



\begin{quote}
Before the new one: say the Malayalam for a signature, then the Malayalam for a duplicate copy.
\end{quote}

\subsection*{You'll want to know: \ml{തിരിച്ചറിയൽരേഖ}}
\textbf{\ml{തിരിച്ചറിയൽരേഖ}} — \emph{tiriccaṟiyalrēkha} — \textquotedblleft{}an identity document\textquotedblright{}.

\begin{quote}
\ml{വിമാനത്താവളത്തിൽ} \ml{തിരിച്ചറിയൽരേഖ} \ml{കാണിക്കണം}.
\end{quote}

\begin{quote}
\emph{vimānattāvaḷattil tiriccaṟiyalrēkha kāṇikkaṇaṁ.}
\end{quote}

\begin{quote}
You must show an identity document at the airport.
\end{quote}

\ml{തിരിച്ചറിയൽ}, recognising, from \ml{തിരിച്ചറിയുക}, plus \ml{രേഖ}, a document.

\begin{cousinweb}
\textbf{\ml{തിരിച്ചറിയൽരേഖ}} (\emph{tiriccaṟiyalrēkha}) is a Malayalam compound built from \ml{തിരിച്ചറിയൽ}, recognition, from the native verb \ml{തിരിച്ചറിയുക} (DEDR 314), and \ml{രേഖ}, a document, from Sanskrit \emph{rekhā}, a line (CDIAL 10810).
\end{cousinweb}

\begin{grammarlens}[title={What you've built}]
One more word for getting around.
\end{grammarlens}

\subsection*{Your turn}
\begin{itemize}
\raggedright
  \item \emph{Say it:} \emph{tiriccaṟiyalrēkha}
  \item \emph{Say it:} \emph{tiriccaṟiyalrēkha}, once more
  \item \emph{Say it:} \emph{vimānattāvaḷattil tiriccaṟiyalrēkha kāṇikkaṇaṁ.} — You must show an identity document at the airport.
  \item \emph{Recall:} say the Malayalam for a signature, then the Malayalam for a duplicate copy, then say \emph{tiriccaṟiyalrēkha} again
\end{itemize}

\begin{tcolorbox}[breakable,colback=teal!4,colframe=teal!35!black,title={Before you move on}]
What does \ml{തിരിച്ചറിയൽരേഖ} mean? (\textquotedblleft{}An identity document\textquotedblright{}.) Say it once more.
\end{tcolorbox}
\section[\texorpdfstring{\ml{കാലാവധി} (kālāvadhi)}{കാലാവധി (kālāvadhi)}]{\ml{കാലാവധി} (kālāvadhi) — a validity period; a term}

\label{lesson:ML-C488-kalavadhi}



\begin{quote}
Before the new one: say the Malayalam for a duplicate copy, then the Malayalam for an identity document.
\end{quote}

\subsection*{You'll want to know: \ml{കാലാവധി}}
\textbf{\ml{കാലാവധി}} — \emph{kālāvadhi} — \textquotedblleft{}a validity period; a term\textquotedblright{}.

\begin{quote}
\ml{എന്റെ} \ml{പാസ്പോർട്ടിന്റെ} \ml{കാലാവധി} \ml{അടുത്ത} \ml{മാസം} \ml{തീരും}.
\end{quote}

\begin{quote}
\emph{enṟe pāspōrṭṭinṟe kālāvadhi aṭutta māsaṁ tīruṁ.}
\end{quote}

\begin{quote}
My passport runs out next month.
\end{quote}

\ml{കാലം}, time, plus \ml{അവധി}, a limit. When it runs out, it \ml{തീരുന്നു}.

\begin{cousinweb}
\textbf{\ml{കാലാവധി}} (\emph{kālāvadhi}) is a Malayalam compound built from two words borrowed from Sanskrit, \emph{kāla}, \textquotedblleft{}time\textquotedblright{}, and \emph{avadhi}, \textquotedblleft{}a limit, a term\textquotedblright{} (Monier-Williams).
\end{cousinweb}

\begin{grammarlens}[title={What you've built}]
One more word for getting around.
\end{grammarlens}

\subsection*{Your turn}
\begin{itemize}
\raggedright
  \item \emph{Say it:} \emph{kālāvadhi}
  \item \emph{Say it:} \emph{kālāvadhi}, once more
  \item \emph{Say it:} \emph{enṟe pāspōrṭṭinṟe kālāvadhi aṭutta māsaṁ tīruṁ.} — My passport runs out next month.
  \item \emph{Recall:} say the Malayalam for a duplicate copy, then the Malayalam for an identity document, then say \emph{kālāvadhi} again
\end{itemize}

\begin{tcolorbox}[breakable,colback=teal!4,colframe=teal!35!black,title={Before you move on}]
What does \ml{കാലാവധി} mean? (\textquotedblleft{}A validity period; a term\textquotedblright{}.) Say it once more.
\end{tcolorbox}
\section[\texorpdfstring{\ml{കാലഹരണപ്പെടുക} (kālaharaṇappeṭuka)}{കാലഹരണപ്പെടുക (kālaharaṇappeṭuka)}]{\ml{കാലഹരണപ്പെടുക} (kālaharaṇappeṭuka) — to expire; to become outdated}

\label{lesson:ML-C488-kalaharanappetuka}



\begin{quote}
Before the new one: say the Malayalam for an identity document, then the Malayalam for a validity period.
\end{quote}

\subsection*{You'll want to know: \ml{കാലഹരണപ്പെടുക}}
\textbf{\ml{കാലഹരണപ്പെടുക}} — \emph{kālaharaṇappeṭuka} — \textquotedblleft{}to expire; to become outdated\textquotedblright{}.

\begin{quote}
\ml{എന്റെ} \ml{ലൈസൻസ്} \ml{കഴിഞ്ഞ} \ml{മാസം} \ml{കാലഹരണപ്പെട്ടു}.
\end{quote}

\begin{quote}
\emph{enṟe laisansŭ kaḻiñña māsaṁ kālaharaṇappeṭṭu.}
\end{quote}

\begin{quote}
My licence expired last month.
\end{quote}

Past: \ml{കാലഹരണപ്പെട്ടു}, \emph{kālaharaṇappeṭṭu}.

A formal verb for papers and medicines; in speech people say the date has passed.

\begin{cousinweb}
\textbf{\ml{കാലഹരണപ്പെടുക}} (\emph{kālaharaṇappeṭuka}) is a Malayalam verb built from Sanskrit \emph{kāla}, \textquotedblleft{}time\textquotedblright{}, and \emph{haraṇa}, \textquotedblleft{}taking away\textquotedblright{} (Monier-Williams), with -\ml{പ്പെടുക}, the ending of verbs that something undergoes.
\end{cousinweb}

\begin{grammarlens}[title={What you've built}]
One more word for getting around.
\end{grammarlens}

\subsection*{Your turn}
\begin{itemize}
\raggedright
  \item \emph{Say it:} \emph{kālaharaṇappeṭuka}
  \item \emph{Say it:} \emph{kālaharaṇappeṭuka}, once more
  \item \emph{Say it:} \emph{enṟe laisansŭ kaḻiñña māsaṁ kālaharaṇappeṭṭu.} — My licence expired last month.
  \item \emph{Recall:} say the Malayalam for an identity document, then the Malayalam for a validity period, then say \emph{kālaharaṇappeṭuka} again
\end{itemize}

\begin{tcolorbox}[breakable,colback=teal!4,colframe=teal!35!black,title={Before you move on}]
What does \ml{കാലഹരണപ്പെടുക} mean? (\textquotedblleft{}To expire; to become outdated\textquotedblright{}.) Say it once more.
\end{tcolorbox}
